% !PRINTED 168
\weberpage{168}

Eine solche Differentialgleichung ist aber in den Formeln des § 23 bereits enthalten.

Es folgt zunächst aus § 23, (14):

\[
d\log \varkappa^2 = i\pi\vartheta_{01}^4\,d\omega,
\]

oder mit Rücksicht auf

\weq{2}{\pi\vartheta_{00}^2 = 2K, \quad \pi\vartheta_{10}^2 = 2\varkappa K, \quad \pi\vartheta_{01}^2 = 2\varkappa'K\ [\S\,42,\ (3),\ (15)]:}

\weq{3}{\pi\,d(\varkappa^2) = 4i\varkappa^2\varkappa'^2K^2\,d\omega,}

und aus § 23, (19):

\[
\frac{d}{d\omega}\frac{dK}{K^2d\omega} = -\frac{4}{\pi^2}\varkappa^2\varkappa'^2K^3.
\]

Führt man vermittelst (3) für $\omega$ die Variable $\varkappa^2$ ein, so folgt

\weq{4}{\frac{d}{d(\varkappa^2)}\left(\varkappa^2\varkappa'^2\frac{dK}{d(\varkappa^2)}\right) = \frac14 K,}

und dies ist die gesuchte Differentialgleichung.

Setzen wir weiter

\weq{5}{-i\omega = \frac{K'}{K},}

\weq{6}{-iK^2d\omega = K\,dK' - K'\,dK,}

so ergibt sich nach (3):

\weq{7}{\varkappa^2\varkappa'^2\left(K\frac{dK'}{d(\varkappa^2)} - K'\frac{dK}{d(\varkappa^2)}\right) = -\frac{\pi}{4},}

also durch abermalige Differentiation:

\[
\frac1K\frac{d}{d(\varkappa^2)}\left(\varkappa^2\varkappa'^2\frac{dK}{d(\varkappa^2)}\right) = \frac{1}{K'}\frac{d}{d(\varkappa^2)}\left(\varkappa^2\varkappa'^2\frac{dK'}{d(\varkappa^2)}\right),
\]

woraus zu ersehen, daß $K'$ das zweite partikulare Integral der Differentialgleichung (4) ist.

% [sic] The source prints the Latin title as “seriem infinitum” (for “seriem infinitam”) and has no comma before Werke.
Diese Differentialgleichung läßt sich durch hypergeometrische Reihen integrieren (\gesperrt{Gauss}, Disq. circa seriem infinitum Werke, Bd. III, S. 125).

Setzen wir für den Augenblick $\varkappa^2 = x$, $K = y$, so lautet die Differentialgleichung (4)

\weq{8}{x(1-x)\frac{d^2y}{dx^2} + (1-2x)\frac{dy}{dx} - \frac14 y = 0,}

und sie geht aus der allgemeinen \gesperrt{Gauss}schen Differentialgleichung

\weq{9}{(x-x^2)\frac{d^2y}{dx^2} + [\gamma-(\alpha+\beta+1)x]\frac{dy}{dx} - \alpha\beta y = 0}
%CONT>>
